\section*{Réalisation cofinale de référence et extinction terminale}

La réalisation qui gouverne ce chapitre est le complément spectral--polaire de l'antécédent enrichi commun. Soit
\[
 \mathcal D_R=C_c^\infty((-R/2,R/2)),\qquad
 \mathcal D_R^{\rm cof}:=\mathcal D_R,\qquad
 \widehat\psi(t)=\int_{\mathbb R}\psi(x)e^{-itx}\,dx ,
\]
et, pour \(\psi,\phi\in\mathcal D_R\),
\[
 h_{\psi,\phi}(z)
 =\widehat\psi(z)\overline{\widehat\phi(\overline z)},\qquad
 \check h_{\psi,\phi}(a)
 =\langle\psi,T_a\phi\rangle ,
 \qquad (T_a\phi)(x)=\phi(x-a).
\]

Soient \(P_+^{\rm sh},P_-^{\rm sh}\) les projecteurs orthogonaux des deux feuillets, et fixons
\begin{equation}\label{eq:amp-rh-owner-q}
 q=\frac59,\qquad
 A=P_+^{\rm sh}+qP_-^{\rm sh},\qquad
 D=\sqrt{1-q^2}\,P_-^{\rm sh},\qquad A^*A+D^*D=I.
\end{equation}
Si \(\Xi_R\) est la restriction de l'antécédent commun, \(Q_R=I-P_R\) son complément orthogonal et \(\jmath_-\) l'inclusion isométrique dans le feuillet antispéculaire, on construit
\begin{equation}\label{eq:amp-rh-owner-binding}
 E_R:=\jmath_-Q_R\Xi_R,\qquad
 AE_R=qE_R,\qquad
 E_R^*E_R=\mathcal B_{{\rm W},R}.
\end{equation}
L'identité de conservation s'itère sans effacer le dernier terme :
\begin{equation}\label{eq:amp-rh-owner-telescopia}
 \mathcal B_{{\rm W},R}
 =\sum_{n=0}^{N-1}(DA^nE_R)^*(DA^nE_R)
 +(A^NE_R)^*(A^NE_R).
\end{equation}
Le terminal est calculé avant le passage à la limite :
\begin{equation}\label{eq:amp-rh-owner-terminal}
 (A^NE_R)^*(A^NE_R)
 =q^{2N}\mathcal B_{{\rm W},R}\longrightarrow0
 \qquad(N\longrightarrow\infty).
\end{equation}
Ce n'est qu'après cette extinction que l'on définit
\begin{equation}\label{eq:amp-rh-owner-l}
 \mathscr L_Rf:=(DA^nE_Rf)_{n\geq0},
 \qquad
 \boxed{\mathcal B_{{\rm W},R}=\mathscr L_R^*\mathscr L_R\succeq0}.
\end{equation}
La compatibilité de restriction exprime ces factorisations comme une forme cofinale unique :
\begin{equation}\label{eq:amp-rh-owner-publicacion}
 B_W=\Xi^*(I-P)\Xi
 =\bigl((I-P)\Xi\bigr)^*\bigl((I-P)\Xi\bigr)\succeq0,
 \qquad
 B_W(f,f)=\|(I-P)\Xi f\|^2.
\end{equation}

\begin{theorem}[Clôture cofinale Riemann--Weil]
\label{thm:amp-rh-cierre-cofinal-propietario}
La forme complète de Guinand--Weil est positive sur la classe cofinale \(\bigcup_{R>0}\mathcal D_R^{\rm cof}=C_c^\infty(\mathbb R)\). Par le critère réversible de Weil, tout zéro non trivial \(\rho\) de la fonction zêta de Riemann satisfait
\[
 \Re\rho=\frac12.
\]
\end{theorem}

\begin{proof}
Les contraintes de \eqref{eq:amp-rh-owner-binding} transforment le terminal de \eqref{eq:amp-rh-owner-telescopia} exactement en le membre de gauche de \eqref{eq:amp-rh-owner-terminal}. Son extinction géométrique permet de passer à \eqref{eq:amp-rh-owner-l} ; la naturalité cofinale produit \eqref{eq:amp-rh-owner-publicacion}. L'identification premier--gamma--polaire avec la fonctionnelle complète et la séparation cofinale du développement principal de référence \eqref{eq:realizaciones-rh-identificacion-gw} permettent d'appliquer le critère de Weil. Les calculs ultérieurs vérifient à nouveau cette expression dans une carte régulée auxiliaire, sans devenir des prémisses du théorème.
\end{proof}

\section*{Contrôle auxiliaire du régulateur commun et identité des trois canaux}

Le calcul régulé qui suit développe la différence de Gram à partir de ses opérateurs constitutifs avant d'en évaluer le signe. Son résidu est un \texttt{CONTRÔLE\_NON\_DÉTERMINANT} : il délimite cette fenêtre particulière et ne remplace ni ne conditionne la réalisation cofinale \eqref{eq:amp-rh-owner-binding}--\eqref{eq:amp-rh-owner-publicacion}. Fixons une fonction \(\vartheta\in C_c^\infty([0,\infty))\) telle que \(0\leq\vartheta\leq1\), \(\vartheta=1\) sur \([0,1]\) et \(\vartheta=0\) sur \([2,\infty)\), et posons
\[
 \vartheta_L(a)=\vartheta(a/L),\qquad \vartheta_L(0)=1 .
\]
C'est l'unique régulateur des longueurs : il s'applique simultanément aux horloges premier--puissance et au canal gamma ; les termes scalaire et polaire sont ses atomes de longueur nulle. Il n'est pas permis de régulariser chaque feuillet avec une coupure différente.

Pour
\[
 \ell_{p,k}=k\log p,\qquad
 w_{p,k}=(\log p)p^{-k/2},\qquad
 \mu_\Gamma(a)=\frac{e^{-a/2}}{1-e^{-2a}},
\]
définissons également
\begin{align*}
 c_\Gamma&=\psi_\Gamma(1/4)-\log\pi<0,\\
 A_\psi&=\widehat\psi(i/2),&
 B_\psi&=\widehat\psi(-i/2),&
 b_\pm(\psi)&=\frac{A_\psi\pm B_\psi}{\sqrt2}.
\end{align*}
Au niveau \(N_m=9^m\), soient \(j_m\) l'inclusion uniforme et \(\widehat\delta_{a,m}\) le canal normalisé de la retenue nonadique. Leurs moments se calculent directement à partir de l'odomètre :
\begin{equation}\label{eq:amp-rh-momentos-odometro-explicitos}
 j_m^*j_m=I,\qquad
 \widehat\delta_{a,m}^*\widehat\delta_{b,m}
 =(I-T_a)^*(I-T_b).
\end{equation}
En particulier, \(\widehat\delta_{a,m}\) dépend de la couture et de \(T_a\), non de la fonction zêta ni de ses zéros.

Les deux colonnes régulées sont
\begin{align}
 \mathcal C^+_{m,R,L}\psi
 ={}&
 \Bigl(
  (\sqrt{\vartheta_L(\ell_{p,k})w_{p,k}}\,
       \widehat\delta_{\ell_{p,k},m}\psi)_{p,k},
 \nonumber\\[-2mm]
 &\hspace{11mm}
  a\longmapsto
  \sqrt{\vartheta_L(a)\mu_\Gamma(a)}\,
       \widehat\delta_{a,m}\psi,\,
  \zeta_+b_+(\psi)
 \Bigr),                                             \label{eq:amp-rh-columna-mas-explicita}\\
 \mathcal C^-_{m,R,L}\psi
 ={}&
 \Bigl(
  (\sqrt{2\vartheta_L(\ell_{p,k})w_{p,k}}\,j_m\psi)_{p,k},\,
  \sqrt{-c_\Gamma}\,j_m\psi,\,
  \zeta_-b_-(\psi)
 \Bigr),                                             \label{eq:amp-rh-columna-menos-explicita}
\end{align}
où les sommes ne parcourent que les longueurs pour lesquelles \(\vartheta_L(\ell_{p,k})\neq0\), et \(\zeta_\pm\) sont les droites polaires orthogonales de la représentation cyclique d'ordre \(12\).

\subsection*{Couplage non décomposable antérieur au signe}

La relation entre \eqref{eq:amp-rh-columna-mas-explicita} et \eqref{eq:amp-rh-columna-menos-explicita} ne s'obtient pas en postulant une contraction entre deux matrices de Gram déjà calculées. Soit
\[
 P_m=\iota_mE_m,\qquad Q_m=I-P_m
\]
l'espérance nonadique explicite. Pour l'horloge unitaire \(\mathscr U_{\ell,m}\), écrivons
\[
 a_{\ell,m}=P_m\mathscr U_{\ell,m}P_m,\qquad
 c_{\ell,m}=Q_m\mathscr U_{\ell,m}P_m ,
\]
de sorte que l'unité des blocs donne
\begin{equation}\label{eq:amp-rh-defecto-schwarz-reloj}
 P_m-a_{\ell,m}^*a_{\ell,m}
 =c_{\ell,m}^*c_{\ell,m}\succeq0.
\end{equation}
Pour \(r=p^{-1/2}\), posons
\[
 s_{p,m}=\frac{N_mr}{1+(N_m-1)r},
 \qquad \xi_{p,m}=\operatorname{arsech}(s_{p,m}),
\]
et soit \(\mathscr S_{\xi_{p,m}}\) la colligation polaire unitaire obtenue par la transformation de Potapov--Ginzburg. Sa composition de Redheffer avec l'horloge vérifie
\begin{align}
 C_{p,m}&=(s_{p,m}I-a_{\log p,m})
          (I-s_{p,m}a_{\log p,m})^{-1},\nonumber\\
 I-C_{p,m}^*C_{p,m}
 &=\bigl(1-s_{p,m}^2\bigr)
   (I-s_{p,m}a_{\log p,m}^*)^{-1}\nonumber\\
 &\quad{}\times c_{\log p,m}^*c_{\log p,m}
   (I-s_{p,m}a_{\log p,m})^{-1}.
   \label{eq:amp-rh-redheffer-defecto}
\end{align}
Le développement de Neumann, avant toute expression spectrale, donne
\begin{equation}\label{eq:amp-rh-torre-prima-explicita}
 \lambda_{p,m}(I-C_{p,m}^*C_{p,m})
 =\sum_{k\geq1}(\log p)p^{-k/2}
  \bigl(2I-T_{k\log p}-T_{k\log p}^*\bigr),
\end{equation}
avec
\[
 \lambda_{p,m}
 =\frac{(\log p)r(1+r)}
 {(1-r)\kappa_{p,m}},
 \qquad
 \kappa_{p,m}
 =\frac{(1-s_{p,m}^2)(N_m-1)}
 {[N_m-(N_m-1)s_{p,m}]^2}.
\]
Ainsi, les puissances d'un nombre premier sont les coefficients d'une unique rétroaction modulaire, non une liste choisie après coup. Soit \(D_{p,L}\) la contraction diagonale du port d'atteignabilité de cette rétroaction, qui multiplie sa coordonnée \(k\) par \(\sqrt{\vartheta_L(k\log p)}\). La régulation ne modifie ni l'horloge ni ses coefficients : elle comprime simultanément les coordonnées de sortie de la même tour.

Le port fini des rôles n'est pas non plus diagonal. Dans les secteurs actifs
\[
 \mathcal H_+
 =\operatorname{span}\{f_3,f_6,f_9\},\qquad
 \mathcal H_-=\operatorname{span}\{f_4,f_8\},
\]
la coisométrie
\begin{equation}\label{eq:amp-rh-acoplamiento-angular}
 C_\angle
 =|f_4\rangle\langle f_3|
  +|f_8\rangle\langle f_9|
\end{equation}
satisfait
\[
 C_\angle C_\angle^*=I_{\mathcal H_-},\qquad
 I_{\mathcal H_+}-C_\angle^*C_\angle=P_6.
\]
La carte isométrique \(U_{W_{24}\to12}\), obtenue à partir des \(5\) octades incidentes et de leur Gram \(4I_5+\frac43J_5\), transporte cette coisométrie à la frontière HMT :
\[
 C_{W_{24}}
 =U_{W_{24}\to12}^*C_\angle U_{W_{24}\to12}.
\]

\begin{theorem}[Contrôle du réseau Paley--Wiener--polaire et de son résidu de port]
\label{thm:amp-rh-acoplamiento-prospectivo}
Pour \(m,R,L\) fixés, définissons
\begin{align}
 \mathscr R^{\rm pr}_{m,L}
 &=
 \mathop{\star}_{\log p\leq2L}
 D_{p,L}\bigl(\mathscr S_{\xi_{p,m}}\star
              \mathscr U_{\log p,m}\bigr),\nonumber\\
 \mathscr R^\Gamma_{m,L}
 &=
 \int_0^{2L}{}^\oplus
 \mathscr U_{a,m}
 \vartheta_L(a)\,d\mu_\Gamma(a),\nonumber\\
 \mathscr C^{\rm TPK}_{m,R,L}
 &=C_{W_{24}}\star
 \bigl(\mathscr R^{\rm pr}_{m,L}
       \oplus\mathscr R^\Gamma_{m,L}\bigr).
 \label{eq:amp-rh-red-prospectiva}
\end{align}
Ce réseau est passif. Dans les espaces énergétiques natifs de ses ports, soient \(J_{+,m,R,L}:=\mathcal C^+_{m,R,L}\), \(J_{-,m,R,L}:=\mathcal C^-_{m,R,L}\) et \(\mathscr C^{\rm src}_{m,R,L}\) le transfert contractif externe de \eqref{eq:amp-rh-red-prospectiva}. On définit
\begin{equation}\label{eq:amp-rh-publicaciones-construidas}
 \begin{aligned}
 \mathcal E_{m,R,L}
 &:=\mathscr C^{\rm src}_{m,R,L}J_{+,m,R,L}
    -J_{-,m,R,L},\\
 \mathscr D_{m,R,L}
 &:=\bigl(I-(\mathscr C^{\rm src}_{m,R,L})^*
                  \mathscr C^{\rm src}_{m,R,L}\bigr)^{1/2},\\
 \mathsf L_{m,R,L}&:=\mathscr D_{m,R,L}J_{+,m,R,L}.
 \end{aligned}
\end{equation}
Le résidu \(\mathcal E_{m,R,L}\) mesure si cette carte auxiliaire de fenêtre reproduit à elle seule l'expression déjà construite par le développement cofinal de référence. Son annulation n'est ni une exigence de la clôture Riemann--Weil ni une prémisse de sa factorisation positive. Aucun coefficient de \eqref{eq:amp-rh-red-prospectiva} ne dépend d'un zéro de \(\zeta\), de la positivité de Weil ou du signe qui sera obtenu.
\end{theorem}

\begin{proof}
Chaque horloge \(\mathscr U_{a,m}\) est unitaire. L'identité \eqref{eq:amp-rh-defecto-schwarz-reloj} identifie son port de retenue. La transformation polaire est unitaire et \eqref{eq:amp-rh-redheffer-defecto} démontre la passivité après fermeture du port interne. La composition de Redheffer d'un nombre fini de colligations passives est encore passive. Le canal gamma s'obtient comme limite de sommes directes dans \([\varepsilon,2L]\) ; l'estimation \(\mu_\Gamma(a)\|(I-T_a)\psi\|^2=O(a)\) lorsque \(a\downarrow0\) donne la limite dans le domaine de forme.

Le développement \eqref{eq:amp-rh-torre-prima-explicita} produit exactement les coefficients premier--puissance. L'intégrale directe produit la seconde entrée de \eqref{eq:amp-rh-columna-mas-explicita}. La transformation de Hadamard au port polaire produit \(b_+\) et \(b_-\), et \eqref{eq:amp-rh-acoplamiento-angular}, conjuguée par \(U_{W_{24}\to12}\), détermine les \(3\) entrées et les \(2\) sorties du même état. Ces opérations construisent le transfert et les deux colonnes natives ; leur différence de sortie est exactement \(\mathcal E_{m,R,L}\) dans \eqref{eq:amp-rh-publicaciones-construidas}. La passivité garantit que \(\mathscr D_{m,R,L}\) est bien définie, mais n'annule pas à elle seule ce résidu auxiliaire ; elle n'a pas non plus à le faire pour la clôture cofinale déjà démontrée.

Ce réseau n'est pas \(C_\angle\otimes I\) : les facteurs \((I-s_{p,m}a_{\log p,m})^{-1}\) mélangent chaque port polaire avec toutes les puissances de l'horloge, et l'intégrale directe mélange le même port avec le continu gamma. Supprimer ces termes modifie la colonne d'atteignabilité et ne constitue pas une simplification de \eqref{eq:amp-rh-red-prospectiva}.
\end{proof}

\subsection*{Développement terme à terme de la différence de Gram}

\begin{lemma}[Identité premier--gamma--polaire avec régulateur commun]
\label{lem:amp-rh-expansion-guinand-weil}
Pour \(\psi,\phi\in\mathcal D_R\),
\begin{equation}\label{eq:amp-rh-diferencia-columnas-regulada}
 \langle J_{+,m,R,L}\psi,J_{+,m,R,L}\phi\rangle
 -\langle J_{-,m,R,L}\psi,J_{-,m,R,L}\phi\rangle
 =
 \mathscr I_{{\rm GW},L}(\psi,\phi),
\end{equation}
où
\begin{align}
 \mathscr I_{{\rm GW},L}(\psi,\phi)
 ={}&
 h_{\psi,\phi}(i/2)+h_{\psi,\phi}(-i/2)
 +c_\Gamma\check h_{\psi,\phi}(0)                 \nonumber\\
 &+\int_0^\infty\vartheta_L(a)\mu_\Gamma(a)
 \bigl(
 2\check h_{\psi,\phi}(0)
 -\check h_{\psi,\phi}(a)
 -\check h_{\psi,\phi}(-a)
 \bigr)\,da                                       \nonumber\\
 &-\sum_{p,k}\vartheta_L(\ell_{p,k})w_{p,k}
 \bigl(
 \check h_{\psi,\phi}(\ell_{p,k})
 +\check h_{\psi,\phi}(-\ell_{p,k})
 \bigr).                              \label{eq:amp-rh-gw-truncado-explicito}
\end{align}
C'est la fonctionnelle complète tronquée de Guinand--Weil : elle contient le terme polaire, le facteur gamma et les \(2\) orientations de chaque puissance de nombre premier sous un régulateur unique.
\end{lemma}

\begin{proof}
Par les définitions \eqref{eq:amp-rh-columna-mas-explicita}--\eqref{eq:amp-rh-columna-menos-explicita},
\begin{align}
 \langle J_+\psi,J_+\phi\rangle
 -\langle J_-\psi,J_-\phi\rangle
 &=
 \langle\mathcal C^+\psi,\mathcal C^+\phi\rangle
 -\langle\mathcal C^-\psi,\mathcal C^-\phi\rangle .
 \label{eq:amp-rh-expansion-inicial}
\end{align}
Dans le canal premier, \eqref{eq:amp-rh-momentos-odometro-explicitos} donne
\begin{align*}
 &\langle(I-T_\ell)\psi,(I-T_\ell)\phi\rangle
 -2\langle\psi,\phi\rangle\\
 &\hspace{18mm}
 =-\check h_{\psi,\phi}(\ell)
  -\check h_{\psi,\phi}(-\ell).
\end{align*}
La multiplication par \(\vartheta_L(\ell_{p,k})w_{p,k}\) suivie de la sommation produit exactement la dernière ligne de \eqref{eq:amp-rh-gw-truncado-explicito}.

Dans le canal gamma,
\[
 \langle(I-T_a)\psi,(I-T_a)\phi\rangle
 =2\check h(0)-\check h(a)-\check h(-a).
\]
En outre, la représentation intégrale de la dérivée logarithmique de gamma donne, pour \(t\in\mathbb R\),
\begin{equation}\label{eq:amp-rh-digamma-longitudes}
 \Re\psi_\Gamma\!\left(\frac14+\frac{it}{2}\right)-\psi_\Gamma(1/4)
 =2\int_0^\infty\mu_\Gamma(a)(1-\cos at)\,da .
\end{equation}
L'intégration contre \(h(t)/(2\pi)\) et l'inversion de Fourier donnent
\begin{align*}
 &-\log\pi\,\check h(0)
 +\frac1{2\pi}\int_{\mathbb R}h(t)
  \Re\psi_\Gamma\!\left(\frac14+\frac{it}{2}\right)\,dt\\
 &\qquad =
 c_\Gamma\check h(0)
 +\int_0^\infty\mu_\Gamma(a)
  [2\check h(0)-\check h(a)-\check h(-a)]\,da .
\end{align*}
La version régulée applique à cette identité la substitution commune
\[
 \mu_\Gamma(a)\longmapsto
 \vartheta_L(a)\mu_\Gamma(a).
\]

Enfin,
\begin{align*}
 \langle b_+(\psi),b_+(\phi)\rangle
 -\langle b_-(\psi),b_-(\phi)\rangle
 &=
 A_\psi\overline{B_\phi}
 +B_\psi\overline{A_\phi}\\
 &=h_{\psi,\phi}(i/2)+h_{\psi,\phi}(-i/2).
\end{align*}
La somme des \(3\) développements est \eqref{eq:amp-rh-gw-truncado-explicito}. Chaque égalité a été calculée avant de connaître le signe de leur somme.
\end{proof}

\section*{Convergence, cofinalité et séparation}

\begin{proposition}[Suppression du régulateur et naturalité cofinale]
\label{prop:amp-rh-convergencia-cofinal}
Pour toute \(\psi,\phi\in\mathcal D_R\), la limite
\[
 \lim_{L\to\infty}\mathscr I_{{\rm GW},L}(\psi,\phi)
 =\mathscr I_{\rm GW}(\psi,\phi)
\]
existe et coïncide avec la forme complète de Guinand--Weil. Si \(R\leq R'\), les inclusions \(\mathcal D_R\hookrightarrow\mathcal D_{R'}\) et les raffinements \(m\mapsto m+1\) entrelacent séparément \(J_{+,m,R,L}\) et \(J_{-,m,R,L}\). Par conséquent, les identités \eqref{eq:amp-rh-diferencia-columnas-regulada} forment un système compatible unique, non une collection de factorisations dépendant de la fenêtre.
\end{proposition}

\begin{proof}
Puisque \(\check h_{\psi,\phi}\in C_c^\infty((-R,R))\), la somme sur les puissances de nombres premiers qui subsiste après annulation des contre-termes scalaires est localement finie : tous les termes avec \(\ell_{p,k}\geq R\) s'annulent. Dans l'intégrale gamma,
\[
 2\check h(0)-\check h(a)-\check h(-a)=O(a^2)
 \quad(a\downarrow0),
\]
tandis que \(\mu_\Gamma(a)\sim(2a)^{-1}\) ; l'intégrande est \(O(a)\). Pour \(a>R\), les deux translations disparaissent et \(\mu_\Gamma(a)=O(e^{-a/2})\). La convergence dominée permet \(L\to\infty\).

L'identité de moments \eqref{eq:amp-rh-momentos-odometro-explicitos} ne dépend pas de \(m\) ; elle définit donc l'isométrie de raffinement \(\widehat\delta_{a,m}\mapsto\widehat\delta_{a,m+1}\), tandis que \(j_m\mapsto j_{m+1}\). Les droites polaires et \(U_{W_{24}\to12}\) restent fixes. Lorsque la fenêtre s'élargit, les nouvelles horloges ont des translations disjointes sur \(\mathcal D_R\) et ajoutent le même terme \(2w_{p,k}\langle\psi,\phi\rangle\) aux deux colonnes avant de s'annuler. Cela démontre les deux entrelacements.
\end{proof}

\paragraph{Contrôle fini hérité de la fibre nonadique.}
Soit \(\mathscr F_{729}\) le facteur local fini à \(729\) états, \(P^{[729]}\) le projecteur sur les \(81\) états grossiers et \(Q^{[729]}=I_{\mathscr F_{729}}-P^{[729]}\). Alors
\begin{equation}\label{eq:amp-rh-729-81-648}
 729=81+648,\qquad
 \operatorname{rank}P^{[729]}=81,\qquad
 \operatorname{rank}Q^{[729]}=648.
\end{equation}
Ce recensement précède le critère de signe et n'identifie pas le complément \(Q^{[729]}\) au résidu de connecteur \(\mathcal E_{m,R,L}\).

\begin{proposition}[Compatibilité cofinale du complément fini]
\label{prop:amp-rh-compatibilidad-complemento-finito}
\label{prop:amp-rh-residual}
Les identités d'expression commune de \eqref{eq:realizaciones-rh-publicaciones-comunes}--\eqref{eq:realizaciones-rh-bw} forment, sous les raffinements \(m\mapsto m+1\) et les inclusions \(\mathcal D_R\hookrightarrow\mathcal D_{R'}\), un système compatible unique. En particulier, la différence des deux matrices de Gram est transportée comme la forme du complément \((I-P)\Xi\), et le recensement \eqref{eq:amp-rh-729-81-648} demeure un contrôle fini de la fibre.
\end{proposition}

\begin{proof}
L'orthogonalité de \(P\) et les expressions communes donnent \eqref{eq:realizaciones-rh-bw}. La \cref{prop:amp-rh-convergencia-cofinal} entrelace séparément les deux colonnes sous raffinement et inclusion ; en les soustrayant, l'identité du complément est conservée dans tout le système cofinal.
\end{proof}

Dans cette carte auxiliaire, on n'affirme pas \(\mathcal E_{m,R,L}=0\) : la \cref{prop:amp-rh-obstruccion-ventana-corta} délimite exclusivement ce modèle régulé. Le recensement retrouvé appartient au complément fini de référence, non au contrôle de fenêtre courte ; le résidu ne rouvre ni ne conditionne donc le théorème cofinal.

\begin{lemma}[La classe cofinale sépare les évaluations spectrales]
\label{lem:amp-rh-separacion-cofinal}
On a l'égalité, et non seulement la densité,
\[
 \bigcup_{R>0}\mathcal D_R=C_c^\infty(\mathbb R).
\]
En outre, pour tous points complexes distincts \(z_1,\ldots,z_n\), l'application
\[
 \mathcal D_R\longrightarrow\mathbb C^n,\qquad
 \psi\longmapsto
 (\widehat\psi(z_1),\ldots,\widehat\psi(z_n))
\]
est surjective pour tout \(R\) suffisamment grand. En particulier, aucune orbite finie d'évaluations du critère de Weil ne reste invisible pour le système cofinal.
\end{lemma}

\begin{proof}
La première affirmation découle de la compacité du support. Pour la seconde, si les fonctionnelles d'évaluation étaient linéairement dépendantes, il existerait des \(c_1,\ldots,c_n\) non tous nuls tels que
\[
 \int_{\mathbb R}\psi(x)
 \left(\sum_{j=1}^nc_je^{-iz_jx}\right)\,dx=0
 \qquad(\psi\in\mathcal D_R).
\]
La fonction analytique \(\sum_jc_je^{-iz_jx}\) s'annulerait sur l'intervalle \((-R/2,R/2)\), donc sur tout \(\mathbb C\). L'indépendance des exponentielles de fréquences distinctes impose \(c_j=0\) pour tout \(j\), d'où une contradiction. Les fonctionnelles sont indépendantes ; une application linéaire vers \(\mathbb C^n\) dont l'adjoint est injectif est surjective.
\end{proof}

\section*{Couplage de 2 parcours et mise à l'épreuve contradictoire du plateau}

Le produit nu \(C_\angle\otimes I\) exigerait
\[
 2\|\psi\|^2\leq\|(I-T_\ell)\psi\|^2
\]
pour toute \(\psi\), inégalité qui échoue sur des plateaux lisses beaucoup plus longs que \(\ell\). Le réseau \eqref{eq:amp-rh-red-prospectiva} n'effectue pas cette identification : il conserve les croisements Paley--Wiener, gamma, premier et polaire avant projection. Le calcul suivant expose cette différence sur un plateau concret.

Soit \(\omega=e^{2\pi i/9}\) et, dans la fibre survivante,
\[
 \kappa_0=u_{81}\otimes\chi_6,\qquad
 \kappa_1=v_{10}\otimes\chi_6,\qquad
 \chi_6(c)=3^{-1}\omega^{6c}.
\]
Ce sont \(2\) parcours orthogonaux du même rôle \(P_6\). Leur expression de frontière produit \(e_0,e_1\) avec
\[
 \langle e_i,\Lambda_{\rm tw}^{\#}e_j\rangle=\delta_{ij}.
\]
Prenons
\[
 a=\frac1{16},\qquad b=2\log2,\qquad
 \psi_w=w^{-1/2}{\bf1}_{[-w/2,w/2]}\quad(w=a,b).
\]
Ces fonctions appartiennent à la fermeture de forme de \(\mathcal D_R\), et leurs régularisations internes appartiennent à \(\mathcal D_R\). Si \(g_{ij}=\mathscr I_{\rm GW}(\psi_{w_i},\psi_{w_j})\), avec \((w_0,w_1)=(a,b)\), le développement précédent donne, pour \(\lambda_n=2n+\tfrac12\),
\begin{align}
 g(w,w)
 ={}&\frac2w\sum_{n\geq0}
 \frac{1-e^{-\lambda_nw}}{\lambda_n^2}
 +c_\Gamma+\frac{32}{w}\sinh^2\!\frac w4 \nonumber\\
 &-2\sum_{k\log p<w}
  w_{p,k}\left(1-\frac{k\log p}{w}\right),
 \label{eq:amp-rh-meseta-diagonal}\\
 g_{01}
 ={}&\frac2{\sqrt{ab}}\sum_{n\geq0}
 \frac{e^{-\lambda_n(b-a)/2}(1-e^{-\lambda_na})}{\lambda_n^2}
 +c_\Gamma\sqrt{\frac ab} \nonumber\\
 &+2A_aA_b-\frac{\log2}{\sqrt2}\sqrt{\frac ab},
 \qquad A_w=\frac{4\sinh(w/4)}{\sqrt w}.
 \label{eq:amp-rh-meseta-cruce}
\end{align}
La majoration des queues géométriques et de la série de \(\lambda_n^{-2}\) produit
\begin{align}
 1.46610954095406&<g_{00}<1.46690960095857,\nonumber\\
 0.06966817455957&<g_{11}<0.06970424464086,\nonumber\\
 -0.008430670493497&<g_{01}<-0.008430670493494,
 \label{eq:amp-rh-meseta-cotas}
\end{align}
et, par multiplication d'intervalles,
\begin{equation}\label{eq:amp-rh-meseta-determinante}
 0.10207009921767<
 \det\begin{pmatrix}g_{00}&g_{01}\\g_{01}&g_{11}\end{pmatrix}
 <0.10217874948627.
\end{equation}
Le Gram complet du plateau et de l'état de base est donc de rang \(2\) ; une seule droite scalaire ne peut l'exprimer.

Posons
\[
 \sigma_b=g_{11}-\frac{g_{01}^2}{g_{00}}>0
\]
et définissons sur \(\mathcal E_2=\operatorname{span}\{\psi_a,\psi_b\}\)
\begin{align}
 H^{(2)}_{m,R}(\alpha\psi_a+\beta\psi_b)
 ={}&
 e_{0,m}\left(
 \sqrt{g_{00}}\,\alpha
 +\frac{g_{01}}{\sqrt{g_{00}}}\,\beta\right)
 +e_{1,m}\sqrt{\sigma_b}\,\beta .
 \label{eq:amp-rh-lector-dos-rutas}
\end{align}
Une multiplication de Cholesky, utilisant \(\langle e_i,\Lambda_m^\#e_j\rangle=\delta_{ij}\), démontre
\begin{equation}\label{eq:amp-rh-meseta-energia-exacta}
 \langle H^{(2)}_{m,R}\psi,
         \Lambda_m^\#H^{(2)}_{m,R}\phi\rangle
 =\mathscr I_{\rm GW}(\psi,\phi)
 \qquad(\psi,\phi\in\mathcal E_2).
\end{equation}
Les coefficients de \eqref{eq:amp-rh-lector-dos-rutas} proviennent des séries explicites \eqref{eq:amp-rh-meseta-diagonal}--\eqref{eq:amp-rh-meseta-cruce} ; aucun zéro n'est introduit et le signe du déterminant n'est pas présupposé : il est démontré dans \eqref{eq:amp-rh-meseta-determinante}. Par continuité des \(3\) canaux, l'identité et la positivité stricte du déterminant persistent pour des régularisations lisses suffisamment proches. C'est l'épreuve contradictoire que le tenseur \(C_\angle\otimes I\) ne satisfait pas, contrairement au couplage non décomposable.

\section*{Résidu exact et délimitation du régulateur}

Définissons
\begin{equation}\label{eq:amp-rh-residual-finito}
 \mathcal R_{m,R,L}
 :=
 (\mathcal C^+_{m,R,L})^*\mathcal C^+_{m,R,L}
 -(\mathcal C^-_{m,R,L})^*\mathcal C^-_{m,R,L}
 -\mathscr I_{{\rm GW},L}.
\end{equation}
Le lemme \ref{lem:amp-rh-expansion-guinand-weil} démontre
\[
 \mathcal R_{m,R,L}=0
\]
terme à terme, et non par définition. Parallèlement, \eqref{eq:amp-rh-publicaciones-construidas} donne, par calcul des normes, le bilan du port
\begin{equation}\label{eq:amp-rh-balance-residual-puerto}
 \boxed{
 \begin{aligned}
 \mathscr I_{{\rm GW},L}(\psi,\psi)
 -\|\mathsf L_{m,R,L}\psi\|^2
 ={}&2\operatorname{Re}
 \langle J_{-,m,R,L}\psi,\mathcal E_{m,R,L}\psi\rangle\\
 &+\|\mathcal E_{m,R,L}\psi\|^2 .
 \end{aligned}}
\end{equation}
En effet, \(\|\mathsf L\psi\|^2
=\|J_+\psi\|^2-\|\mathscr C^{\rm src}J_+\psi\|^2\) et \(\mathscr C^{\rm src}J_+=J_-+\mathcal E\). Ainsi, dans cette carte auxiliaire, une annulation démontrée de \(\mathcal E\) produirait le carré positif
\begin{equation}\label{eq:amp-rh-cuadrado-regulado}
 \mathcal E_{m,R,L}=0
 \quad\Longrightarrow\quad
 \mathscr I_{{\rm GW},L}(\psi,\psi)
 =\|\mathsf L_{m,R,L}\psi\|^2\geq0.
\end{equation}
L'implication n'autorise pas à supposer la prémisse et ne fait pas partie de la chaîne de référence de la clôture cofinale.

\begin{proposition}[Obstruction de fenêtre courte]
\label{prop:amp-rh-obstruccion-ventana-corta}
Pour les colonnes \eqref{eq:amp-rh-columna-mas-explicita}--\eqref{eq:amp-rh-columna-menos-explicita}, le résidu \(\mathcal E_{m,R,L}\) ne peut s'annuler pour toute \(\psi\in\mathcal D_R\) et tout \(L>0\).
\end{proposition}

\begin{proof}
Choisissons
\[
 0\ne\psi\in\mathcal D_R,\qquad
 \widehat\psi(i/2)=\widehat\psi(-i/2)=0.
\]
Une telle fonction existe puisque deux fonctionnelles linéaires sont imposées sur l'espace de dimension infinie \(\mathcal D_R\). Si \(2L<\log2\), la somme sur les puissances de nombres premiers de \eqref{eq:amp-rh-gw-truncado-explicito} est vide et le canal polaire s'annule. Pour \(h_{\psi,\psi}\), on a \(\check h_{\psi,\psi}(0)=\|\psi\|_2^2\), de sorte que
\begin{equation}\label{eq:amp-rh-ventana-corta-expansion}
 \mathscr I_{{\rm GW},L}(\psi,\psi)
 =c_\Gamma\|\psi\|_2^2+
 \int_0^{2L}\vartheta_L(a)\mu_\Gamma(a)
       \|(I-T_a)\psi\|_2^2\,da .
\end{equation}
Puisque \(\|(I-T_a)\psi\|_2\leq a\|\psi'\|_2\) et \(\mu_\Gamma(a)=O(a^{-1})\) lorsque \(a\downarrow0\), l'intégrale de \eqref{eq:amp-rh-ventana-corta-expansion} est \(O_\psi(L^2)\). Comme \(c_\Gamma<0\), pour \(L\) suffisamment petit, on obtient
\[
 \mathscr I_{{\rm GW},L}(\psi,\psi)<0.
\]
Si \(\mathcal E_{m,R,L}\psi=0\), l'implication \eqref{eq:amp-rh-cuadrado-regulado} donnerait le signe opposé. Le résidu ne peut donc s'annuler avec le quantificateur indiqué.
\end{proof}

L'égalité \(\mathcal R_{m,R,L}=0\) identifie correctement la différence des deux matrices de Gram à la fonctionnelle tronquée ; elle ne la transforme pas en carré positif. Ce connecteur régulé est donc un \texttt{CONTRÔLE\_NON\_DÉTERMINANT} : il n'intervient pas dans la chaîne probatoire \eqref{eq:amp-rh-owner-binding}--\eqref{eq:amp-rh-owner-publicacion}. La proposition \ref{prop:amp-rh-obstruccion-ventana-corta} délimite exclusivement les formules \eqref{eq:amp-rh-columna-mas-explicita}--\eqref{eq:amp-rh-gw-truncado-explicito} ; son résultat ne dégrade ni la positivité cofinale ni le \cref{thm:amp-rh-cierre-cofinal-propietario}.
